\section{模型评价与推广}

\subsection{各问结论的衔接}

四个问题的结论沿同一条证据链衔接。问题一给出可复现的相对质量评分与 1M 规模上的配比响应面，
但响应面的损失水平不能跨规模迁移，质量评分也只覆盖 17 个配比领域中的 6 个；问题二因此只能采纳
不含质量的经典标度律，包含质量的候选形式作为未通过预设检验的结果保留，经典标度律在独立观测
数据上的中位相对误差约为 8\%；问题三在这一标度律的有效域内求解，最优配置随预算依次经过内点最优、
数据量取上界、预算有剩余三个区间，最高代表性预算中只有很小一部分能在有效域内使用，质量只能以
基线计入；问题四在经典标度律的损失尺度上识别不出损失到得分的斜率，因此不经由标度律换算，而直接
以历史得分给出 12 个月的前沿延续基线 50.89，并把算力增长放缓表述为给定转移假设下的敏感性结果。
前一问没有识别的关系，后一问都没有当作已知条件使用；各问的否定性结论与肯定性结论一样，界定了
本文结论的适用范围。

\subsection{模型的优点}

\begin{enumerate}
  \item \textbf{证据分级贯穿全文。}直接证据、模型条件推论、半合成校准、外推与敏感性、诊断与
        未决始终分开表述；在拟合前检验主数据的性质，避免把生成公式的恢复误认为新规律的发现。
  \item \textbf{判据预先设定。}配比模型的选择规则、广义标度律的采纳条件、分析总体与时间口径
        都在查看结果之前固定，否定性结论（广义形式不采纳、跨规模不迁移、映射无斜率）同样作为
        结果报告。
  \item \textbf{约束与转移有明确的数学定义。}算力约束保留为不等式并限定在有效域内求解，区分预算
        用尽与预算有剩余；结构性转移定义为约束激活状态的改变，并以预算弹性的跳变检验。
  \item \textbf{不确定性按数据结构度量。}质量评分按文档自助、标度律按训练轨迹重抽样、
        前沿按整月块自助，并把时间口径、总体定义等方法选择的影响单独列出。
\end{enumerate}

\subsection{模型的不足}

\begin{enumerate}
  \item 质量评分是相对代理量，冲突成因的解释来自领域分布模式，未对具体指标对逐一检验。
  \item 配比响应面只在 1M 规模上可作绝对预测，现有数据不足以支持向 10B/70B 的外推。
  \item 质量未能以可检验的方式进入标度律；经典标度律的经验支持只来自两张较小的独立观测表。
  \item 资源配置以经典标度律为前提，而该标度律在主拟合数据上只证明了生成公式可以恢复；上下文
        长度只有 5 个观测取值；质量基线 $Q_0$ 没有数值，非基线的质量配置未能计算。
  \item 前沿预测依赖 9 个月的密集历史和线性趋势假设，时间口径的影响大于模型本身的不确定性；
        算力放缓情景依赖从预训练子集转移而来的占比。
\end{enumerate}

\subsection{模型的推广}

体系平衡的质量聚合与冲突度量适用于任何多打分器的数据筛选场景；在拟合前用已发表公式检验
数据是否由公式生成的做法，可推广到一切使用组织方或第三方汇编数据的建模任务；有效域内的不等式
约束配置与以约束激活状态定义的结构性转移，适用于其他只在有限范围内得到识别的资源配置模型，一旦
给定 $Q_0$ 的数值和质量进入损失的可检验形式，非基线的质量配置也可以在同一框架中计算；把以参数量与
以训练算力为尺度的分解并列，可用于评估其他领域中规模扩张的贡献。
